\color{unidarkgreen}{\subsection{Внешнее отключение (ВО)}}
\color{black}

\begin{enumerate}[label=\arabic{section}.\arabic{subsection}.\arabic*,  align=left, labelsep=4pt, leftmargin=0pt, itemindent=57pt]

\item
Внешнее отключение предназначено для приема внешних сигналов и формирования требуемого
действия на отключение выключателей всех сторон автотрансформатора. 

\begin{figure}[H]
\centering
\includegraphics[width=0.8\textwidth,height=0.8\textheight,keepaspectratio]{\fbpath/041.1761 ВО/img/VO.pdf}
\captionof{figure}{Функциональная схема ВО-1}\label{pic_vo:VO}
\end{figure}

\item
В составе логики внешнего отключения предусмотрен прием трех сигналов внешнего отключения \textbf{«Внеш.
откл1(2,3)»} с возможность контроля по току. ВО-1(2,3) вводится в работу программным переключателем \textbf{«ВО-
1(2,3)»}. Контроль по току определяется программным переключателем \textbf{«Конт I ВО-1(2,3)»}, а порог
срабатывания тока задается уставкой \textbf{«Iср ВО-1(2,3)»}. Действие внешнего отключения осуществляется с
регулируемыми независимыми выдержками времени \textbf{«Тср ВО-1(2,3)»}.

\item
Режим действия внешнего отключения задается с помощью программного переключателя \textbf{«Режим ВО-1(2,3)»}:
\begin{itemize}
\item \textit{«Без АПВ»} -- Внешнее отключение с запретом АПВ;
\item \textit{«С АПВ»} -- Внешнее отключение без запрета АПВ.
\end{itemize} 

\item
Оперативный ввод/вывод внешнего отключения производится с помощью БУ (блоков управления функциями) \textit{«Вывод ВО»}. 

\small
\begin{longtable}{|>{\centering\arraybackslash}m{5.3cm}|>{\centering\arraybackslash}m{4.2cm}|>{\centering\arraybackslash}m{1cm}|>{\centering\arraybackslash}m{3.3cm}|>{\centering\arraybackslash}m{1.8cm}|}
% Первая страница: полный заголовок
\caption{Перечень уставочных параметров ВО\hfill\vspace{-0.5\baselineskip}}\label{vo:tbl1}\\
\hline
\rowcolor{gray!30}
Наименование уставки & Диапазон & Шаг & Значение по умолчанию & Единица измерения \\
\hline
\endfirsthead

% Промежуточные страницы: «Продолжение таблицы»
\caption*{\hspace{3pt}\emph{Продолжение таблицы \ref{vo:tbl1}}\hfill\ }\vspace{-0.5\baselineskip} \\
\hline
\rowcolor{gray!30}
Наименование уставки & Диапазон & Шаг & Значение по умолчанию & Единица измерения \\
\hline
\endhead

%===m>ВО|-
%===m
\end{longtable}
\normalsize	

\end{enumerate}